\chapter{Soteriologie in der frühen Kirche}
Die Ansichten zu einer bestimmten Doktrin in der Kirchengeschichte sind nicht unser primärer Weg,
Gottes Sicht auf die Doktrin zu verstehen, aber können uns doch helfen, einige Ansichten als späte Entwicklung auszusortieren.
Es ist mein Standpunkt und meine Argumentation in diesem Kapitel, dass PSA wie in \ref{psa-definition} definiert keine Sicht war, die in den ersten Jahrhunderten der Kirchengeschichte existiert.
Im folgenden werden wir einige Zitate ins Auge nehmen, die \Cite{psa-in-church-history}
verwendet um zu dem Schluss zu kommen ``that the first thousand years of ancient Christianity frequently espoused the teaching that Jesus suffered death, punishment, and a curse for fallen humanity as the penalty for human sin''\Footcite[199, Abstract]{psa-in-church-history}.
Gleichzeitig werde ich - wie in den Kommentaren zu biblischem Material - positiv beschreiben, was diese Väter der Kirche tatsächlich über Jesus' erlösendes Wirken denken.

Bevor wir zu den einzelnen Texten übergehen möchte ich darauf hinweisen, dass PSA für viele seiner Befürworter `der Herz des Evangeliums'\Footcite{psa-is-the-heart-of-the-gospel} ist.
Was erwarten wir in den christlichen Werken der Kirche, wenn das der Fall ist?
Warum lehnen sowohl die Katholische als auch die Orthodoxe Kirche PSA ab, wenn diese Ansicht nicht nur ein Bestandteil, sondern \textit{der} Bestandteil der Christlichen Botschaft ist?
Wenn PSA das Herz des Evangeliums ist, dann erwarten wir in den Schriften der frühen Kirche eine klare Formulierung dieser Doktrin.
Was wir stattdessen finden sollen die nächsten Seiten demonstrieren.

Zunächst findet sich in den Konzilen keine Diskussion und dementsprechend keine Aussagen über die Versöhnungslehre an sich. Gehen wir jetzt also über zu den Aussagen einzelner Kirchenväter
.\footnote{Nicht alle Zitate aus \Cite{psa-in-church-history} sind es überhaupt Wert, diskutiert zu werden.
Als Beispiel nenne ich hier \Cite[49:6]{clement-epistula-ad-corinthios}: ``In Liebe hat der Herr uns angenommen; wegen der Liebe, die er zu uns trug, hat unser Herr Jesus Christus sein Blut hingegeben für uns nach Gottes Willen, sein Fleisch für unser Fleisch, seine Seele für unsere Seelen.''
Nichts in dieser Aussage ist in irgend einer Form distinktiv für PSA.
Christus hat sein Blut und sein Fleisch für uns hingegeben.
Diese Aussage sagt nichts dazu, ob der Vater den Sohn für unsere Sünden bestraft hat.}

\section{Athanasius}
Ich stelle die Diskussion über Athanasius voran, weil sein Verständnis sowohl vom Problem, dass Jesus löst als auch der Lösung, die er dafür gibt leicht verständlich in einem einzelnen Werk dargelegt ist.
Gleichzeitig wird ein gutes Verständnis von Athanasius' Ansicht refokusieren, mit welcher Frage wir tatsächlich an die Texte der Kirchenväter herantreten müssen um zu beantworten, ob sie PSA lehren oder nicht.

In \Cite{de-incarnatione-verbi} gibt Athanasius eine detaillierte Exposition und Verteidigung der Natur der Inkarnation des Wortes, und seines Wirkens.
In diesem Werk können wir also auch erwarten, dass das Erlösungswerk Gottes im Fleischgewordenen Wort zur Sprache kommt und wir werden nicht enttäuscht.
Angefangen von der Erschaffung der Welt erklärt Athanasius, welchen Effekt Sünde hatte, und was das zu lösende Problem seiner Meinung nach ist:

\begin{quote}
Gott hat uns ja nicht nur aus dem Nichts erschaffen, sondern durch die Gnade des Logos
uns auch ein gottgleiches Leben verliehen.
Da aber die Menschen das Ewige von sich wiesen und auf Eingebung des Teufels sich dem Vergänglichen zuwandten, so haben sie selber
damit ihre Vergänglichkeit im Tode verschuldet; sie waren ja wohl, wie vorhin bemerkt,
von Natur aus sterblich, wären aber diesem natürlichen Los dank ihrer Teilnahme am Logos entronnen, wenn sie gut geblieben wären.
Denn um des Logos willen, der mit ihnen
war, hätte die naturgemäße Vernichtung sie nicht betroffen, wie auch die Weisheit sagt:
„Gott hat den Menschen zur Unverweslichkeit erschaffen und als Bild seiner eigenen Ewigkeit;
aber durch den Neid des Teufels kam der Tod in die Welt“\bibfoot{Wisdom}{2}{24}
Wie aber das geschehen war, begann das Sterben der Menschen;
die Vergänglichkeit gewann nunmehr und entfaltete um so stärkere Macht über das ganze Geschlecht denn der natürliche Urstand, als
sie auch die Drohung Gottes wegen der Übertretung des Gebotes ihnen gegenüber voraus hatte.
\attrib{\Cite[5]{de-incarnatione-verbi}}
\end{quote}

Für die Soteriologie in \Cite{de-incarnatione-verbi} ist das zu lösende Problem,
dass sich der Mensch durch seine `Übertretung des Gebotes' von dem natürlichen Zustand
(Unsterblichkeit wegen des geschaffen Seins im Leben des Logos) der Vergänglichkeit (dem Tod) zugewandt hat.

Wir können Athanasius juristische Sprache an folgendem Abschnitt sehr direkt nachvollziehen:
\begin{quote}
So also hat Gott den Menschen erschaffen und ihn in der Unsterblichkeit belassen wollen.
Die Menschen jedoch würdigten den geistigen Verkehr mit Gott wenig, kehrten sich
davon ab, erdachten und ersannen sich die Bosheit, wie im ersten Teil ausgeführt wurde,
und verfielen dem angedrohten Todesurteil. Jetzt sollten sie auch nicht mehr so bleiben,
wie sie geschaffen worden sind, vielmehr sanken sie entsprechend ihrer Denkart immer
tiefer, und der Tod wurde ihr Gewaltherr. Denn die Übertretung des Gebotes warf sie auf
ihren natürlichen Urstand zurück, so daß sie, wie aus dem Nichts geworden, so auch mit
Recht nach Ablauf der Zeit den Verlust ihrer Existenz zu gewärtigen hatten.
Denn wenn
es in ihrer Natur lag, einmal nicht zu sein, und sie erst durch das Eingreifen und die Menschenliebe des Logos ins Dasein gerufen wurden,
so ergab sich als natürliche Folge, daß die Menschen mit dem Verlust ihrer Gottesvorstellung und mit ihrer Abkehr zum Nichtseienden
- nichtseiend ist das Böse, seiend das Gute, weil ja vom Seienden Gott ausgegangen -
auch ihrer ewigen Existenz verlustig gingen, das heißt aber, daß sie der Auflösung anheimfielen und im Tod und in der Verwesung verblieben.
Tatsächlich ist ja der Mensch von
Natur aus sterblich, da er aus dem Nichts entstanden ist. Doch dank seiner Ähnlichkeit
mit dem Seienden hätte er in dem Falle, daß er sie mit einer richtigen Herzensstellung zu
ihm bewahrt hätte, die naturgemäße Auflösung von sich ferngehalten und wäre unverweslich geblieben, wie ja die Weisheit sagt: „Die Beobachtung der Gebote ist die Sicherung der
Unverweslichkeit“\bibfoot{Wisdom}{6}{19}.
\attrib{\Cite[4]{de-incarnatione-verbi}}
\end{quote}

Die Natürliche Ordnung ist, dass der Mensch abseits seiner `richtigen Herzensstellung zu Gott' stirbt.
In Athanasius' Worten ist, das `angedrohte Todesurteil' über die Sünde des Menschen: (a) Unterwerfung unter die Herrschaft des Todes (b) Rückkehr zum natürlichen Zustand abseits des Seienden Logos - Sterblichkeit.
An dieser Stelle wird unmissverständlich klar: für Athanasius ist der Tod ein Feind Gottes, der ihm Diametral gegenübersteht\footnote{nichtseiend ist das Böse, seiend das Gute, weil ja vom Seienden Gott ausgegangen}.
Durch die Erschaffung aus dem Nichts hat das Logos den Menschen Lebendig gemacht.
Durch seine Abwendung von Gott hat der Mensch Tod auf sich gebracht, und zwar \textit{als Folge der natürlichen Ordnung}.
Tod in Athanasius ist keine retributive Strafe Gottes.
Tod ist eine natürlich folgende Konsequenz, die Athanasius im gleichen Atemzug ohne Probleme das `angedrohte Todesurteil' nennen kann.
Wir sehen also, dass das Problem der Tod ist - und zwar Tod als natürliche Ordnung des Menschen abseits der lebensspendenden Einheit mit dem Logos.
Das Problem ist weder Gottes Zorn über Sünde der besänftigt werden muss, noch ein Problem der ausstehenden Strafe Gottes für Sünde, die der Mensch nicht tragen kann.
Das Problem ist, dass der Mensch abseits von Gott Tod ist.
Gleichzeitig ist natürlich genau diese Ordnung von Gott gegeben und der Tod des Menschen, die Übergabe des Menschen an das Nichtssein das aus Sünde resultiert, Gottes Gesetz.
\footnote{So schreibt Athanasius wenig später
``Denn, wie vorhin gesagt, der Tod gewann jetzt über uns gesetzliche Gewalt, und es war unmöglich, dem Gesetze zu entrinnen, weil es von Gott wegen der Sünde erlassen worden war.''
(\Cite[6]{de-incarnatione-verbi}).
Es ist nicht wahr zu sagen, dass Athanasius keine juristische Sprache verwendet, bloß ist nicht Gottes juristische Forderung das Problem, sondern `der Tod gewann die gesetzliche Gewalt':
Der Tod, in Perversion des gerechten Urteil Gottes (Menschen abseits von ihm sterben), kann Sünde verwenden um Tod zu erzeugen.
Die Relationale Trennung des Menschen zu Gott führt zu einer Ontologischen Trennung vom göttlichen Leben des Logos.
Siehe auch \Cite{deification-in-athanasius} für eine deutlich detailliertere Diskussion über die primären Themen in Athanasius' Soteriologie.}

In der Folge macht Athanasius sehr deutlich, warum Buße allein nicht ausreicht, um dieses Problem zu lösen.
(a) Gott wäre nicht wahrhaftig gewesen in der Aufstellung seines Gesetzes, dass Sünde zu Tod führt
(b) Reue hebt den Natürlichen Zustand (also genau dieses gottgegebene, natürliche Gesetz an sich) nicht auf, sondern beendet nur das Sündigen.
Athanasius schreibt ``Wenn also nur Sünde, und nicht auch Vernichtung als deren Folge dagewesen wäre, dann wäre es mit der Reue gut gewesen.''\Footcite[7]{de-incarnatione-verbi}.
In anderen Worten: Wenn nur das Sündigen an sich das Problem wäre, wäre Reue genug um ihr ein Ende zu setzen.
Aber `Vernichtung als deren Folge' (nie in Athanasius als retributive Strafe) kann nicht allein durch Buße beendet werden. Es braucht eine bessere Heilung, die nur das inkarnierte Logos erwirken kann:

\begin{quote}
Er sah das vernünftige Geschlecht zugrunde gehen und den Tod mit seiner Verwesung herrschen über die Menschen;
er sah, wie auch die Strafandrohung für die Sünde uns im Banne des Verderbens festhalte und eine Befreiung daraus vor der Erfüllung des Gesetzes unangebracht wäre;
er sah auch das Unziemliche, das im Falle des Unterganges der Wesen, deren Schöpfer er selber war, gelegen wäre;
er sah auch die überflutende Bosheit der Menschen und wie sie diese nachgerade bis zur Unerträglichkeit zu ihrem eigenen Verderben steigerten;
er sah endlich alle Menschen als die Beute des Todes.
Deshalb erbarmte er sich unseres Geschlechtes, hatte Mitleid mit unserer
Schwachheit, ließ sich herab zu unserer Vergänglichkeit, duldete die Herrschaft des Todes
nicht, und um die Schöpfung gegen den Tod zu schützen und das Werk seines Vaters an
den Menschen nicht vergeblich sein zu lassen, nahm er einen Leib an, und zwar keinen
anderen als den unsrigen.
\attrib{\Cite[8]{de-incarnatione-verbi}}
\end{quote}
Die Probleme sind, in Reihenfolge
(a) die Herrschaft des Todes
(b) die Strafandrohung für Sünde\footnote{
Wie mehrfach gesagt: Die Strafe/das Urteil ist in Athanasius, dass Gott den Menschen seinem natürlichen Nichtssein abseits des lebensspendenden Logos überlässt; die Strafe ist keine retributive Ausschüttung Gottes Zorns über die Sünde.}
- Tod - kann nur aufgehoben werden, wenn das Gesetz erfüllt wird\footnote{Denn ansonsten wäre Gottes Ehre und Wahrhaftigkeit verletzt. Siehe dazu Fußnote 21 in der Zitierten Version.}
(c) Gott will seine eigenen Geschöpfe nicht wieder dem Nichtssein preisgeben
(d) die Bosheit der Menschen führt zu ihrem eigenen Verderben
(e) die Herrschaft des Todes.
Dementsprechend ist auch die Antwort geformt (a) als herabsteigen des Logos zu unserer Vergänglichkeit (b) widerstand gegen und Schutz vor dem Tod (c) damit das Werk des Vaters nicht vergeblich ist.
``Und das tat er aus Liebe zu den Menschen, damit alle in ihm sterben
und so das Gesetz von der Verwesung der Menschen aufgehoben würde, da ja seine Macht
am Leibe des Herrn sich erschöpft hat und bei den gleichartigen Menschen keinen Zugang
mehr finden kann. Auch wollte er die Menschen, die in die Verweslichkeit zurückgefallen
waren, wieder zur Unverweslichkeit erheben und sie vom Tode zu neuem Leben erwecken,
indem er durch die Aneignung des Leibes und die Gnade der Auferstehung den Tod in
ihnen wie eine Stoppel im Feuer vernichtete.''\Footcite[8]{de-incarnatione-verbi}.
Das Gesetz (dass Menschen Sterben) wird aufgehoben, indem sein Tod die `Macht des Gesetzes' erchöpft wird.
Alle sterben in ihm, der Tod kann das Leben des Logos nicht überwinden, und so wird der Tod vernichtet.

\begin{quote}
	Daher hat er den Leib, den er angenommen,
als eine Weihegabe und als ganz makelloses Schlachtopfer in den Tod gegeben und verscheuchte alsobald von allen seinesgleichen den Tod durch
das stellvertretende Opfer.
Denn weil erhaben über alle, hat natürlich der Logos Gottes
mit der Hingabe seines Tempels und der leiblichen Werkstätte ein Entgelt für das Leben
aller entrichtet und die Schuld in seinem Tod bezahlt.
\attrib{\Cite[9]{de-incarnatione-verbi}}
\end{quote}
Somit kommen wir jetzt mit einem klaren Verständnis von Athanasius Problemstellung zu diesen Sätzen, die in \Cite{psa-in-church-history} als Beweis dafür verwendet werden, dass Athanasius PSA lehrt.
Ja, in Athanasius ist Jesus als stellvertretendes Opfer gestorben.
\footnote{τῇ προσφορᾷ τοῦ καταλλήλου NPNF(\Cite{de-incarnatione-verbi-npnf}) übersetzt `by the offering of an equivalent'.
Der Fokus liegt offensichtlicherweise auf der Gleichheit des Opfers, geschaffen durch die Inkarnation.
Die Distinktion zwischen Substitution (Jesus stirbt, damit die Menschheit leben kann) und Repräsentation/Partizipation (alle sterben in ihm, damit das Gesetzt von der Verwesung aufgehoben wird) ist keine, an die Athanasius denkt.
}
Ja, in Athanasius hat Jesus durch seinen Tod die Strafe für die Sünde getragen.
Nein, in keiner Art und Weise lehrt Athanasius PSA.
Wir müssen jeden Author für sich sprechen lassen - in seiner Sprache, über seine Probleme.
Das suchen nach kurzen Zitaten, die unsere Position verteidigen ist außerordentlich gefährlich, führt zu Verwirrung und wir müssen den Denkern der Vergangenheit mit mehr Respekt begegnen.

\Cite{psa-in-church-history} zitiert weiter aus Athanasius' erster Rede gegen die Arianer:
\begin{quote}
	Denn damals wurde die Welt wie ein Schuldiger vom Gesetze gerichtet. Nun aber nahm das Wort das Gericht auf sich, litt im Leibe für alle und schenkte allen das Heil.
	\attrib{\Cite[I.60.]{orationes-contra-arianos}}
\end{quote}
Bevor wir im folgenden betrachten, was Athanasius hier meint, ist direkt offensichtlich, dass dieser Satz nicht verwendet werden kann um PSA zu verteidigen. Es ist in diesem Zitat das Gesetz, dass richtet, nicht der Vater.
Das Gericht des Gesetzes ist, was das Wort auf sich nimmt.
Was also ist das Problem, dass die Inkarnation hier löst?
Im vorhergehenden Absatz lesen wir:
\begin{quote}
Das Gesetz aber wurde von den Engeln verkündet und hat niemand vollkommen gemacht,
da es der Ankunft des Wortes bedurfte, wie Paulus gesagt hat.
Die Ankunft des Wortes aber hat des Vaters Werk vollendet. Und damals herrschte von Adam bis auf Moses der Tod,
das Erscheinen des Wortes aber machte dem Tode ein Ende.
Und wir sterben nicht mehr alle in Adam, sondern werden alle in Christus lebendig gemacht.
Und damals wurde von Dan bis Bersabee das Gesetz verkündigt, und in Judäa allein war Gott bekannt.
Jetzt aber ist über die ganze Erde ihr Ruf gedrungen.
\attrib{\Cite[I.59]{orationes-contra-arianos}}
\end{quote}
In diesem Abschnitt seiner Rede will Athanasius begründen, dass das inkarnierte Wort besser ist als die Engel.
Er vergleicht also das durch die Engel verkündete Gesetzt\bibfoot{Hebrews}{2}{2} mit der Erlösung im inkarnierten Wort.
Was ist es, das in Athanasius' Argument fehlt, und was bewirkt die Inkarnation des Wortes?
Seit Adam herrscht der Tod, und durch das Wort wird dem Tod ein Ende gemacht.
Wieder ist das Problem der Tod, und die Antwort die Inkarnation.
Wir finden auch in dieser Ausführung keinen Bezug dazu, dass Gottes Zorn am Kreuz auf Jesus ausgeschüttet wurde.
Wir finden keine Aussage dazu, dass Jesus für unsere Schuld bestraft wurde.
Auch Vergebung als Konzept spielt absolut keine Rolle.
Was wir stattdessen sehen ist ganz einfach: Das Wort nahm das Gericht (des Gesetzes, das ist der Tod) auf sich, und schenkte allen das Heil (das Leben).
Keine dieser Aussagen ist distinktiv für PSA, und keine der distinktiven Teile von PSA (der Vater bestraft den Sohn für unsere Vergehen) taucht in Athanasius' denken auf.

Wir halten nach dieser längeren Diskussion über Athanasius Soteriologie die folgenden Punkte fest:
(a) wir müssen alle Vorsicht walten lassen, wenn wir die Schriften anderer Kulturen und Zeiten lesen; das selbe Wort kann in einem anderen Kontext eine gänzlich andere Bedeutung haben
(b) die Aussage, die wir tatsächlich in den Kirchenvätern suchen um zu sagen, dass sie PSA verteidigen ist diese: `der Vater hat den Sohn stellvertretend für die Sünde der Menschen bestraft; auf diese Art und Weise bewirkt das Kreuz Vergebung'.
Diese Aussage ist die minimalistische Form von PSA, und diese Aussage werden wir auch im weiteren Verlauf nicht einmal in der Kirche der Antike finden.

\section{Brief an Diognet}
Der Brief \Cite{mathetes-epistula-ad-diognetus} ist laut einigen Verteidigern von PSA `the single best description of penal substitution in the first few centuries, and quite possibly in the history of the church'\Footcite{arnold-psa-in-church-history}, und `if I didn’t know better you could have told me Calvin or Luther wrote this'\Footcite{diognetus-on-christs-atonement}.
Und wenn wir die Stelle in Frage zitieren scheint hier tatsächlich eine wahre Perle der PSA, versteckt im zweiten Nachchristlichen Jahrhundert, zu finden zu sein:
\begin{quote}
Als aber das Mass unserer Ungerechtigkeit voll und es völlig klar geworden war, dass als ihr Lohn Strafe und Tod uns erwarte,
und als der Zeitpunkt gekommen war, den Gott vorausbestimmt hatte, um fortan seine Güte und Macht zu offenbaren,
- o überschwengliche Menschenfreundlichkeit und Liebe Gottes! -
da hasste und verstiess er uns nicht und gedachte nicht des Bösen, sondern war langmütig und geduldig und nahm aus Erbarmen selbst unsere Sünden auf sich;
er selbst gab den eigenen Sohn als Lösepreis für uns, den Heiligen für die Unheiligen,
den Unschuldigen für die Sünder, den Gerechten für die Ungerechten, den Unvergänglichen für die Vergänglichen, den Unsterblichen für die Sterblichen.
Denn was anders war imstande, unsere Sünden zu verdecken als seine Gerechtigkeit?
In wem konnten wir Missetäter und Gottlose gerechtfertigt werden, wenn nicht allein im Sohne Gottes?
Welch süsser Tausch, welch unerforschliches Walten, welch unverhoffte Wohltat,
dass die Ungerechtigkeit vieler in einem Gerechten verborgen würde und die Gerechtigkeit eines einzigen viele Sünder rechtfertige!
\attrib{\Cite[9]{mathetes-epistula-ad-diognetus}}
\end{quote}
Hallelujah! Aber ist das PSA?
In dieser Stelle erscheint das Problem durchaus strafrechtlichen Charakter zu haben (Lohn der Sünde ist Strafe und Tod).
Direkt vor der zitierten Stelle steht weiter, dass wir aus uns selbst nicht in das Reich Gottes gelangen können, aber durch die Kraft Gottes dazu befähigt werden würden.
Hier ist also ganz deutlich ein Problem zu sehen, dass der Mensch aus seiner Natur (seiner Sünde) gerettet werden muss. Wie geschieht diese Erlösung?
`[Gott] nahm aus Erbarmen selbst unsere Sünde auf sich, er selbst gab den eigenen Sohn als Lösepreis für uns', also ist es in diesem Brief \textit{nicht} der Sohn, der 
die Sünden auf sich nahm, denn das selbe Subjekt gibt dann den eigenen Sohn als Lösepreis.
Dieses `die Sünden auf sich nehmen' kann also nicht als Imputation von Schuld verstanden werden, weil dann die Sünde dem Vater imputiert wird.
Stattdessen kann das griechische αὐτὸς τὰς ἡμετέρας ἁμαρτίας ἀνεδέξατο genauso unproblematisch bedeuten
`er nahm die Sünden auf sich: es hat ihn etwas gekostet, sich um sie zu kümmern'\footnote{
	Vergleiche \Cite[ἀναδέχομαι]{lsj}. Die einfache Bedeutung des Wortes `auf sich nehmen' zum Wohle anderer (ohne rechtliche Imputation) ist in dieser Stelle völlig unproblematisch.
	Eine tatsächliches `annehmen der Schuld für ein Vergehen' ist ebenfalls in der semantischen Reichweite dieses Wortes, aber in unserem Fall durch Grammatik und Theologie ausgeschlossen.
} und diese Interpretation ist hier absolut zu bevorzugen, um nicht den Vater als Objekt einer Imputation von Schuld zu verstehen.
Christus ist hier ein Lösepreis für uns - etwas, dass der Vater `hingegeben hat obwohl es etwas kostet', um uns zu befreien.
Er gab `den Heiligen für die Unheiligen', nicht `substitutionell an unserer Stelle', sondern `um uns damit zu nutzen' - `für uns', nicht `an unserer Stelle'.
Auch der tatsächliche Effekt, den der Author hier beschreibt ist nicht, dass die Strafe Gottes für unsere Sünde stattdessen auf Jesus ausgegossen wurde, und so unsere Schuld abgegolten ist.
Der Effekt ist ganz im Gegenteil, unsere Sünden zu verdecken in seiner Gerechtigkeit.

Das Leben und Wirken und Sterben des Sohnes in Gehorsam zum Vater ist die Grundlage, `in dem wir gerechtfertigt werden'.
Die Sprache ist (a) von Gerechtfertigung durch Jesus und (b) der Rettung des Menschen, der aus sich heraus nicht in der Lage ist, zum Leben zu gelangen.
Das Ausgießen der Strafe Gottes für unsere Sünden auf den Sohn ist kein Bestandteil dieser Passage,
sondern die Antwort auf unsere Schuld ist `nichtanrechnen' und `verdecken'.\footnote{
Ich sage nicht, dass PSA mit diesem Text völlig unvereinbar wäre, es ist bloß nicht in dieser Passage zu finden.}

\section{Johannes Chrysostomus}
Mit Johannes Chrysostomus kommen wir jetzt zu der Passage aus den frühen Kirchenvätern, die einer modernen Form von PSA mit Abstand am nächsten kommt.
So lesen wir in seiner Predigt zu \bibverse{IICorinthians}{5}{21}:
\begin{quote}
Denken wir uns, ein König sieht einen Menschen, einen Räuber und Missethäter, den eben die Strafe ereilt,
und der König gibt seinen geliebten, eingebornen, ebenbildlichen Sohn in den Tod und überträgt nebst dem Tode auch die Schuld von jenem Menschen auf den Sohn,
der doch keine Schuld auf sich hat, und befreit so den Verurtheilten nicht bloß von der Strafe,
sondern auch von der Schmach, und erhebt ihn dann noch zu großer Herrschaft: [...]
\attrib{\Cite[11.IV.]{chrysostomus-2-cor}}
\end{quote}
Hier verwendet Chrysostomus eine Analogie, um zu erklären, wie schrecklich es ist nach der Vergebung durch Jesus Wirken weiter zu sündigen.
Seine Analogie ist die, dass der König (Gott der Vater) den Tod und auch die Schuld jedes Menschen auf Jesus legt, um den Menschen zu erhöhen.
Diese Passage klingt tatsächlich stark an PSA an, aber PSA als komplette Theorie ist mit Chrysostomus' gesamtem theologischen Weltbild nicht vereinbar, wie wir an seiner Behandlung von Vers 19 sehen können:
\begin{quote}
	Denn nicht die Thatsache allein ist wunderbar, daß Gott Freund geworden, sondern auch die Weise, wie er es geworden. Und welches ist diese Weise? Indem er ihnen die Sünden erließ; denn anders war es nicht möglich. Darum fährt Paulus fort: „Indem er ihnen nicht anrechnete ihre Übertretungen.“ Hätte uns Gott für die Versündigungen zur Rechenschaft ziehen wollen, so waren wir alle verloren; denn Alle waren gestorben. Aber trotz der Zahl und Größe der Sünden hat Gott statt der Bestrafung sogar sich versöhnt; er hat die Sünden nicht bloß erlassen, sondern gar nicht angerechnet. So müssen denn auch wir den Feinden vergeben, damit wir ebenfalls der gleichen Vergebung theilhaftig werden.
	\attrib{\Cite[11.II]{chrysostomus-2-cor}}
\end{quote}
Gott hat `keine Rechenschaft verlangt', er hat `nicht bestraft', er hat `sie gar nicht angerechnet'.
Alle diese Aussagen stehen im direkten Widerspruch zu PSA, denn dort kann Gott Sünde nicht `nicht anrechnen'.
Die einzige Möglichkeit, damit Gott in PSA Sünden erlassen kann ist, indem er Sünden stattdessen Jesus anrechnet.

In seinem Kommentar zu \bibverse{Gal}{3}{13} lesen wir von Chrysostomus folgendes:
\begin{quote}
Nun war aber doch das Volk jenem anderen Fluche verfallen: „Verflucht sei jeglicher, welcher nicht beharrt bei allem, was geschrieben steht in dem Buche dieses Gesetzes.“
Wie reimt sich das zusammen?
Nun, das Volk war dem Fluche verfallen; denn es beharrte nicht, und es gab keinen, der das ganze Gesetz erfüllt hätte.
Christus aber tauschte sich für diesen Fluch jenen anderen ein, der da sagt: „Verflucht sei jeglicher, welcher am Holze hängt.“
Da nämlich verflucht ist, wer am Holze hängt, und verflucht auch, wer das Gesetz übertritt, so durfte, wer letzteren Fluch lösen wollte, ihm nicht selber verfallen, mußte aber trotzdem einen Fluch auf sich nehmen; deshalb nahm er jenen für diesen auf sich und hat so durch den einen den anderen gelöst.
Und gleichwie für einen zum Tode Verurteilten ein anderer schuldlos den Tod annimmt und ihn der Strafe entreißt, also hat auch Christus getan.
Denn da er nicht dem Fluche (aus) der Übertretung unterworfen war, nahm er statt dessen den Fluch des Kreuzes auf sich, um jenen Fluch zu lösen.
\attrib{\Cite[III.3.]{chrysostomus-gal}}
\end{quote}
Chrysostomus ist hier sehr klar darin, was er unter `er wurde Fluch für uns' versteht.
In seinem denken übernimmt Jesus \textit{nicht} den Fluch des Gesetztes, unter dem der Mensch ist auf eine substitutionelle Weise.
Stattdessen übernimmt er einen \textit{anderen} Fluch, nämlich den des am-Kreuz-hängens\footnote{Die selbe Sicht vertritt \Cite[237f]{paul-and-the-law-in-chrysostom}}.
Auf diese Art und weise stirbt er also zum Wohle des zu Tode Verurteilten, und entreißt ihn von der Strafe (dem Tod): nicht indem er stattdessen dieselbe Strafe auf sich nimmt,
sondern indem er ihn durch das Aufnehmen eines anderen Fluches repräsentiert/austauscht.\footnote{
Ich denke nicht, dass man aus dieser Stelle herauslesen kann, wie diese Repräsentation/Substitution, beziehungsweise das `lösen des Fluches des Gesetzes' für Chrysostomus mechanisch geschieht.
Das Zentrum des Erlösungswerkes in Chrysostomus' Denken ist ganz eindeutig Christus Victor, wie seine Passah-Homilie zeigt.}

Der folgende Text stammt aus einer Homilie zu Himmelfahrt, und hierin finden wir nach meiner Einschätzung den Text in der frühen Kirche, der PSA am nächsten kommt.
Leider ist dieser Text nicht in deutscher Übersetzung vorhanden, weswegen ich ihn hier auf Englisch zitiere.
\begin{quote}
	See then what happens: the mediator is the Son of him who makes consolation, not a human, nor an angel, nor an archangel, nor any servant.
	And what does the mediator do? The work of the mediator.
	Just as when two people turn their backs on each other and do not wish to be reconciled, someone else must come to intervene and break down the enmity between them, so this is what Christ did.
	God was angry with us, we had turned away from God, the Master who loves mankind, and by putting himself in between, Christ reconciled both natures. 

And how did he come between them? He took on himself the punishment that we deserved from the Father and endured the disgrace and insults that we inflicted on God.
Do you desire to learn how he assumed both that punishment from on high and these insults here below? It is said, Christ redeemed us from the curse of the law by becoming a curse for us\bibfoot{Galatians}{3}{13}.
See then how he received the punishment inflicted from above? Consider how he also endured the insults inflicted here below.
The insults of those who insult you, he says, have fallen on me\bibfoot{Psalms}{68}{10}.
See how he dissolved the enmity, how he did not cease to do, suffer, and painstakingly perform all things until he had brought the adversary and enemy back to God himself and made him a friend?
\attrib{\Cite{chrysostomus-fourth-panegyric}}
\end{quote}
Die parallelen zu PSA sind stark, aber nicht exakt: Zunächst ist das Problem tatsächlich Gottes Zorn auf die Menschen, insbesondere sein gerechtes Gericht über Menschliche Sünde.
Jesus tritt explizit in eine Mittlerrolle zwischen einem `angry God' und `turned away humanity'.
Als dieser Mittler nahm er auf sich `the punishment that we deserved from the Father'.
Wir müssen trotzdem vorsichtig damit umgehen, was Chrysostomus hier sagt und was er nicht sagt, denn die Art und Weise, auf die Jesus `the punishment' auf sich nimmt ist explizit so, wie in \bibverse{Gal}{3}{13} geschrieben.
Aus seiner Auslegung dieser selben Stelle wissen wir, wie Chrysostomus die Übernahme des Fluches bzw. der Strafe versteht:
Jesus \textit{starb}, die Strafe die wir gerecht verdient haben. Er war bereit zu sterben, aber nicht indem er \textit{denselben Fluch}, also \textit{die Strafe als Bestrafung} übernommen hat, sondern \textit{einen anderen Fluch}, nämlich \textit{den Tod als äquivalent der Strafe} trug.
In anderen Worten: Chrysostomus sagt nicht, dass Jesus von Gott stellvertretend für uns bestraft wurde.
Er sagt, dass er bereit war, zu sterben - etwas, das wir verdient hätten, nicht er - um Menschen mit Gott zu versöhnen.
Der kleine aber sehr signifikante Unterschied zwischen Chrysostomus' Aussage und PSA ist gerade genau der Kern von PSA:
Hat Jesus diese Versöhnung bewirkt, \textit{indem er von Gott bestraft wurde}?
Chrysostomus' Homilie zu Himmelfahrt \Cite{chrysostomus-fourth-panegyric} beantwortet dieses Frage, wenn wir Chrysostomus Verständnis von \bibverse{Galatians}{3}{13} mit verwenden nicht mit einem Ja.
Um zu demonstrieren, dass mein Verständnis an dieser Stelle keine unnötige Haarspalterei ist sehen wir uns als nächstes Chrysostomus' Kommentar zu \bibverse{Colossians}{2}{13} an:
\begin{quote}
	Also; Wir waren samt und sonders der Sünde und Strafe verfallen; er nahm selbst die Strafe auf sich und hob dadurch Sünde und Strafe auf; die Strafe aber erlitt er am Kreuze.
	Dort nun heftete er den Schuldbrief an; sodann zerriß er ihn, „wie einer, der Macht hat“.
	— Was für einen Schuldbrief? Entweder er meint damit das, was die Israeliten zu Moses sprachen:
	„Alles, was Gott gesagt hat, wollen wir tun und befolgen“; oder wenn nicht dies, daß wir Gott Gehorsam schuldig sind; oder wenn auch das nicht, daß der Teufel den Schuldbrief in Händen hatte, welchen Gott gegen Adam ausstellte, als er sprach:
	„An welchem Tage du von dem Baume issest, wirst du sterben.“
	Diesen Schuldbrief also hatte der Teufel in Händen; und Christus gab ihn uns nicht zurück, sondern riß ihn selbst entzwei zum Zeichen, daß er die Schuld mit Freuden erlasse.
	\attrib{\Cite[6.3]{chrysostomus-col}}
\end{quote}
Auch hier finden wir eine Aussage, die fast wie PSA klingt, aber in Wirklichkeit etwas ganz anderes bedeutet, denn es ist \textit{der Teufel}, der den Schuldschein in der Hand hält.
Das Gesetz ist von Gott gegeben und gut, wie Chrysostomus' immer wieder betont.
Durch unsere Sünden aber führt das Gesetz zu einer Schuld, und in Chrysostomus' Denken ist es \textit{der Teufel, nicht Gott}, der diesen Schuldschein argwöhnisch über unseren Köpfen hält.
Christus nimmt die Strafe (den Tod) auf sich; zerreißt den Schuldschein, und erlässt uns die Schuld mit Freuden.
Nirgendwo finden wir hier die Idee, dass der Vater den Schuldschein über unserem Kopf hält und Christus die gerechte Strafe des Vaters befriedigt, damit der Vater vergeben kann.
Der Teufel ist der Feind, seine Waffe ist die Strafe des Gesetztes, Christus besiegt den Tod und erlässt die Schuld.
Genau darüber beginnt Chrysostomus dann zu reden: ``Während er sich nämlich Hoffnung machte, Christus selbst in seine Gewalt zu bekommen, büßte er sogar alle diejenigen ein, deren er sich schon bemächtigt hatte: als der Leib des Gottmenschen ans Kreuz geschlagen wurde, da standen die Toten auf.''\Footcite[6.3]{chrysostomus-col}

Wieder ein anderes Bild für Jesus Werk am Kreuz finden wir in seiner Homilie zu \bibverse{Romans}{8}{3}:
\begin{quote}
	Solange die Sünde an den Menschen Sünder fand, brachte sie mit Fug und Recht den Tod über sie; als sie aber einmal auf einen Körper traf, der ohne Sünde war, und ihn doch dem Tode überlieferte, da war sie selbst schuldbar geworden, und das Todesurteil wurde über sie gesprochen.
	[...]
	Siehst du da, wie überall die Sünde verurteilt wird, nicht das Fleisch\footnote{Hier: Jesus' menschlicher Körper}? Wie dieses als Sieger erscheint und als Richter über die Sünde
	\attrib{\Cite[14.5.]{chrysostomus-rom}}
\end{quote}
Wieder sehen wir hier das Bild: Christus siegt durch seinen Tod, indem er die Sünde verurteilte (nicht, indem Gott ihn stellvertretend verurteilt).
Die Sünde, der Teufel, der Tod als personifizierter Agent ist das Problem. Christus siegt über sie und bringt Gericht über sie.\footnote{
Diese Idee finden wir auch in anderen Vätern als sogenanntes `bait-and-switch-Modell' (Lockvogelmodell), auch wenn es in verschiedenen Formen ausformuliert ist.}

Zusammenfassend sehen wir, dass der Zorn Gottes über den Menschen in Chrysostomus sehr real ist, und ein echtes Problem.
Jesus tritt als Mediator auf zwischen der sündlosen, heiligen Natur Gottes und der gefallenen, sündigen Natur des Menschen.
Er ist bereit, dafür den Tod (die Strafe des Gesetzes) stellvertretend für Menschen auf sich zu nehmen, \textit{aber nicht indem er statt uns bestraft wurde},
sondern indem er durch den Tod mit dem Menschen eins wird und so die Naturen zusammenführen kann und vor allem:
damit er als Mensch über Sünde und Tod Siegen kann, die die Menschen unterjocht hatten.
Gottes Vergebung aber ist auch in Chrysostomus nicht auf Gottes Todesurteil über Christus als stellvertretendes Opfer prädikiert.

\section{Cyrill}
tl;dr:
Cyrill ist wichtig, weil seine Katechetischen Briefe genau das sind - eine Zusammenfassung des christlichen Glaubens, den neue Konvertiten zu lernen haben.
PSA spielt darin keine Rolle, was ein Problem für die Position darstellt, dass PSA `der Kern des Evangeliums' ist.

% - Cyrill - Catechetical letters

\section{John of Damascus}
tl;dr:
Von Johannes von Damaskus haben wir eine art systematische Theologie, in der PSA wiederum keine Rolle spielt.
% - John of Damascus - The exact exposition of the orthodox faith

\section{Gregory of Nazianzus}
tl;dr:
Gregor von Nazianz hält eine Sicht, die sich als bestes als `Wiederhergestellt Ikone' zusammenfassen lässt.
Jesus heilt unsere korrupte Natur, um uns vollständiger als Imago Dei hinzustellen.

% - Gregory of Naziansus
% 	- Contra Apollinarius Ep.Cl.
% 		- Soteriology = Christology
% 	- Oration 30
% 	- Oration 37
% 	- Oration 39

\section{Eusebeius}
% - Proof of the gospel

\section{Gregory of Nyssa}
% - Gregory of Nyssa

\section{Augustinus}
% 

% kann ich das streichen?
% ich hab keine Lust, cur deus homo zu lesen
\chapter{Penal Substitutionary Atonement in der Kirchengeschichte}
TODO: ggf. fällt diese Kapitel vollständig weg?
\section{Anselm}
tl;dr:
Selbst in Cur Deus Homo finden wir keine exakte Darstellung von PSA, sondern eher eine Satisfaktionstheorie.
% - Cur Deus Homo
\section{Die Reformation}
tl;dr:
Einige Zitate von den Reformatoren die zeigen, wie hier PSA, etwa in seiner modernen Form wirklich auftritt.

